% LaTeX companion of hw04.typ. Both formats are editable.
\section{Homework 4 - submitted work}

\subsection{Exercise 28 - inverse of a linear transformation}

For

\(T\begin{pmatrix}
x_{1} \\
x_{2} \\
x_{3} \\
x_{4}
\end{pmatrix} = \begin{pmatrix}
22 & 13 & 8 & 3 \\
{- 16} & {- 3} & {- 2} & {- 2} \\
8 & 9 & 7 & 2 \\
5 & 4 & 3 & 1
\end{pmatrix}\begin{pmatrix}
x_{1} \\
x_{2} \\
x_{3} \\
x_{4}
\end{pmatrix},\)

the standard matrix is

\(A = \begin{pmatrix}
22 & 13 & 8 & 3 \\
{- 16} & {- 3} & {- 2} & {- 2} \\
8 & 9 & 7 & 2 \\
5 & 4 & 3 & 1
\end{pmatrix}.\)

The submitted elimination of \(\left( A\  \middle| \ I_{4} \right)\) ends at

\(\begin{pmatrix}
I_{4} & | & \begin{pmatrix}
1 & {- 2} & 9 & {- 25} \\
{- 2} & 5 & {- 22} & 60 \\
4 & {- 9} & 41 & {- 112} \\
{- 9} & 17 & {- 80} & 222
\end{pmatrix}
\end{pmatrix}.\)

Thus

\(A^{- 1} = \begin{pmatrix}
1 & {- 2} & 9 & {- 25} \\
{- 2} & 5 & {- 22} & 60 \\
4 & {- 9} & 41 & {- 112} \\
{- 9} & 17 & {- 80} & 222
\end{pmatrix},\)

and \(T^{- 1}(x) = A^{- 1}x\) for all \(x \in \mathbb{R}^{4}\).

\subsection{Exercise 30 - invertibility}

Let \(A = \begin{pmatrix}
0 & 1 & b \\
{- 1} & 0 & c \\
{- b} & {- c} & 0
\end{pmatrix}\). By Theorem 2.4.3, \(A\) is invertible exactly when \(\text{rref}(A) = I_{3}\). Row reduction gives

\(\begin{pmatrix}
0 & 1 & b \\
{- 1} & 0 & c \\
{- b} & {- c} & 0
\end{pmatrix}\rightarrow\begin{pmatrix}
1 & 0 & {- c} \\
0 & 1 & b \\
0 & 0 & 0
\end{pmatrix}.\)

Therefore \(A\) is not invertible, regardless of the values of \(b,c\).

\subsection{Exercise 42 - permutation matrices}

By elementary transformations that change row order, any permutation matrix can be transformed into \(I_{n}\). Hence its rref is \(I_{n}\), so it is invertible by Theorem 2.4.3. If \(A\) is an arbitrary permutation matrix, then

\(\text{rref}\left( A\  \middle| \ I_{n} \right) = \left( I_{n}\  \middle| \ B \right)\)

and \(B\) is the inverse of \(A\). Since the calculation only changes row order, every row of \(B\) is chosen from \(I_{n}\) without repetition. Thus \(B\) is again a permutation matrix.

\subsection{Exercise 6 - kernel}

For \(A = \begin{pmatrix}
1 & 1 & 1 \\
1 & 2 & 3
\end{pmatrix}\), the submitted row reduction is

\(\begin{pmatrix}
1 & 1 & 1 & | & 0 \\
1 & 2 & 3 & | & 0
\end{pmatrix}\rightarrow\begin{pmatrix}
1 & 0 & {- 1} & | & 0 \\
0 & 1 & 2 & | & 0
\end{pmatrix}.\)

Hence

\(\begin{pmatrix}
x_{1} \\
x_{2} \\
x_{3}
\end{pmatrix} = \begin{pmatrix}
t \\
{- 2t} \\
t
\end{pmatrix} = t\begin{pmatrix}
1 \\
{- 2} \\
1
\end{pmatrix},\quad t \in \mathbb{R},\)

so \(\text{ker}(A) = \text{span}\left( \begin{pmatrix}
1 \\
{- 2} \\
1
\end{pmatrix} \right)\).

\subsection{Exercise 14 - image}

For \(A = \begin{pmatrix}
1 & 2 & 3 \\
1 & 2 & 3 \\
1 & 2 & 3
\end{pmatrix}\),

\(T(x) = Ax = \begin{pmatrix}
1 \\
1 \\
1
\end{pmatrix}x_{1} + \begin{pmatrix}
2 \\
2 \\
2
\end{pmatrix}x_{2} + \begin{pmatrix}
3 \\
3 \\
3
\end{pmatrix}x_{3} = \begin{pmatrix}
1 \\
1 \\
1
\end{pmatrix}\left( {x_{1} + 2x_{2} + 3x_{3}} \right).\)

Therefore \(\text{im}(A) = \text{span}\left( \begin{pmatrix}
1 \\
1 \\
1
\end{pmatrix} \right)\).

\section{Part B}

\subsection{Problem 1 - nilpotent transformations}

Let \(T:\mathbb{R}^{n}\rightarrow\mathbb{R}^{n}\) have standard matrix \(A\).

\subsubsection{(a)}

We prove by induction on \(k\) that the standard matrix of \(T^{k}\) is \(A^{k}\). For \(k = 1\), \(T(x) = Ax = A^{1}x\). Assume \(T^{n{(x)}} = A^{n}x\). Using the composition rule and associativity,

\(T^{n + 1}(x) = T\left( T^{n{(x)}} \right) = T\left( {A^{n}x} \right) = A\left( {A^{n}x} \right) = A^{n + 1}x.\)

Thus the statement holds for all \(k\).

\subsubsection{(b)}

Assume that \(T\) is nilpotent. Then, for some positive integer \(k\), \(T^{k{(x)}} = 0\) for all \(x \in \mathbb{R}^{n}\). By part (a), \(A^{k}x = 0\) for all \(x\). Suppose, for contradiction, that \(A\) is invertible. Then \(A^{k}\) is invertible, with inverse \(\left( {A^{- 1} \cdot A^{- 1} \cdot \ldots \cdot A^{- 1}} \right)\) (\(k\) factors). Thus \(\text{rank}\left( A^{k} \right) = n\), its rref is \(I_{n}\), and the augmented system \(\left( A^{k}\  \middle| \ 0 \right)\) cannot have a solution for every right-hand side as required. This contradicts \(A^{k}x = 0\) for all \(x\). Therefore \(A\) is not invertible.

\subsubsection{(c)}

If \(k = 1\), then \(T = 0\), hence \(A = 0\) and \(A - I_{n} = - I_{n}\) is invertible. For \(k \geq 2\), set

\(B = I_{n} + A + A^{2} + \ldots + A^{k - 1}.\)

Then

\(\left( {I_{n} - A} \right)B = I_{n} - A^{k} = I_{n}\)

and, similarly, \(B\left( {I_{n} - A} \right) = I_{n}\). Thus \(I_{n} - A\) is invertible, and so is \(A - I_{n} = - \left( {I_{n} - A} \right)\).

\subsection{\texorpdfstring{Problem 2 - the vector space \(\mathcal{F}\left( {S,V} \right)\)}{Problem 2 - the vector space \textbackslash mathcal\{F\}\textbackslash left( \{S,V\} \textbackslash right)}}

For \(f,g,h \in \mathcal{F}\left( {S,V} \right)\) and \(x \in S\), pointwise addition gives the vector-space axioms:

\[\left( {f + g} \right)(x) \in V,\quad\left( {f + g} \right)(x) + h(x) = f(x) + \left( {g(x) + h(x)} \right),\quad f(x) + g(x) = g(x) + f(x),\]

and the function \(p(x) = 0_{V}\) is an additive identity. For each \(f\), the function \(n(x) = - f(x)\) is its additive inverse. For scalars \(\alpha,\beta\),

\[\alpha\left( {f + g} \right)(x) = \alpha f(x) + \alpha g(x),\quad\left( {\alpha + \beta} \right)f(x) = \alpha f(x) + \beta f(x),\]

\(\alpha\left( {\beta f(x)} \right) = \left( {\alpha\beta} \right)f(x),\quad 1f(x) = f(x).\)

All expressions are defined in \(V\), so \(\mathcal{F}\left( {S,V} \right)\) is a vector space.

The element \(0_{\mathcal{F}{({S,V})}}\) is the function \(x\mapsto 0_{V}\), whereas \(0_{V}\) is an element of \(V\); they are different elements although the function value is always \(0_{V}\).

\(\mathcal{F}\left( {V,S} \right)\) is not necessarily a vector space. For example, take \(S = \left\{ {1,2} \right\}\) with \(1 + 2 = 3\), which is not in \(S\). The constant functions \(f(v) = 1\) and \(g(v) = 2\) lie in \(\mathcal{F}\left( {V,S} \right)\), but \(f + g\) would have value \(3\), hence is not a function into \(S\).

Finally,

\(\mathcal{P} \subset \mathcal{F}\left( {\mathbb{R},\mathbb{R}} \right),\quad\mathcal{P}_{n} \subset \mathcal{F}\left( {\mathbb{R},\mathbb{R}} \right),\quad C^{\infty} \subset \mathcal{F}\left( {\mathbb{R},\mathbb{R}} \right),\)

where \(\mathcal{P}\) is the set of all real polynomial functions, \(\mathcal{P}_{n}\) those of degree at most \(n\), and \(C^{\infty}\) the infinitely differentiable real functions.

\subsection{Problem 3 - polynomial transformation}

Let \(T(p)(t) = p'(t) + p(0)\).

\subsubsection{(a)}

For \(p_{1},p_{2} \in \mathcal{P}\),

\(T\left( {p_{1} + p_{2}} \right) = \left( {p_{1} + p_{2}} \right)'(t) + \left( {p_{1} + p_{2}} \right)(0) = T\left( p_{1} \right) + T\left( p_{2} \right),\)

and for \(k \in \mathbb{R}\),

\(T\left( {kp} \right)(t) = kp'(t) + kp(0) = kT(p)(t).\)

Thus \(T\) is linear.

\subsubsection{(b)}

\(T_{n}:\mathcal{P}_{n}\rightarrow\mathcal{P}_{n}\) is not surjective: \(p(t) = t^{n}\) is in the target, but no element of \(\mathcal{P}_{n}\) maps to it because differentiation lowers the degree by one. It is not injective: for

\(p_{1}(t) = 1 + 2t + t^{2},\quad p_{2}(t) = 2 + t + t^{2},\)

we have \(T_{n{(p_{1})}} = 3 + 2t = T_{n{(p_{2})}}\) although \(p_{1} \neq p_{2}\).

\subsubsection{(c)}

\(T\) is not injective by the same counterexample. It is surjective: if

\(p(t) = k + a_{1}t + a_{2}t^{2} + \ldots + a_{m}t^{m},\)

take

\(q(t) = kt + \frac{1}{2}a_{1}t^{2} + \frac{1}{3}a_{2}t^{3} + \ldots + \frac{1}{m + 1}a_{m}t^{m + 1}.\)

Then \(q'(t) + q(0) = p(t)\).

\subsection{Problem 4 - left multiplication}

For \(L_{A}:\mathbb{R}^{n \times n}\rightarrow\mathbb{R}^{n \times n}\), \(L_{A{(B)}} = AB\):

\subsubsection{(a)}

For \(B,C \in \mathbb{R}^{n \times n}\) and \(k \in \mathbb{R}\),

\(L_{A{({B + C})}} = A\left( {B + C} \right) = AB + AC = L_{A{(B)}} + L_{A{(C)}},\)

\(L_{A{({kB})}} = A\left( {kB} \right) = kAB = kL_{A{(B)}}.\)

So \(L_{A}\) is linear.

\subsubsection{(b)}

If \(A\) is invertible and \(L_{A{(B_{1})}} = L_{A{(B_{2})}}\), multiplication by \(A^{- 1}\) gives \(B_{1} = B_{2}\), so \(L_{A}\) is injective and hence invertible. Conversely, if \(L_{A}\) is invertible, it is surjective. Thus some \(C\) satisfies \(L_{A{(C)}} = I_{n}\), so \(AC = I_{n}\) and \(A\) is invertible.

\subsubsection{(c)}

If \(L_{A} = L_{B}\), then evaluating both maps at \(I_{n}\) gives \(A = L_{A{(I_{n})}} = L_{B{(I_{n})}} = B\). Hence \(A\mapsto L_{A}\) is injective.

\subsubsection{(d)}

The map \(L:\mathbb{R}^{n \times n}\rightarrow\mathcal{F}\) is not surjective. The constant function \(f(C) = I_{n}\) is in \(\mathcal{F}\), but if \(L(A) = f\), then \(AC = I_{n}\) for all \(C\). Taking \(C = 0_{n \times n}\) is impossible.

\subsection{Problem 5 - rotations and projection}

For \(T = \text{Rot}_{- 80{^\circ}} \circ \text{Proj}_{y} \circ \text{Rot}_{35{^\circ}}\):

\subsubsection{(a)}

\(\text{Rot}_{35{^\circ}}\) has image \(\mathbb{R}^{2}\). Projection onto the \(y\)-axis has image the \(y\)-axis, and rotation clockwise by \(80{^\circ}\) sends this to a line \(80{^\circ}\) clockwise from the \(y\)-axis. Thus the angle between \(\text{im}(T)\) and the \(x\)-axis is \(90{^\circ} - 80{^\circ} = 10{^\circ}\).

\subsubsection{(b)}

The final rotation does not affect the kernel. Projection kills exactly the \(x\)-axis, and undoing the initial \(35{^\circ}\) counter-clockwise rotation gives a kernel line \(35{^\circ}\) counter-clockwise from the \(x\)-axis.

\subsubsection{(c)}

For \(T_{\varphi,\theta} = \text{Rot}_{\varphi} \circ \text{Proj}_{y} \circ \text{Rot}_{\theta}\), the image is a line \(- \varphi\) from the \(y\)-axis and the kernel is a line \(\theta\) from the \(x\)-axis. They agree when

\(\varphi = \frac{\pi}{2} - \theta + k\pi,\)

equivalently \(\theta - \varphi = \frac{\pi}{2} + k\pi\) for \(k \in \mathbb{Z}\).
